
Current code generation benchmarks measure whether agents produce correct code, not how efficiently they debug failing code. An agent passing 95\% of tests may have made 10 strategic root-cause fixes or 100 random modifications---existing metrics cannot distinguish these scenarios. This work introduces fix-impact-ratio, a metric that discriminates strategic debugging (ratio 3.90: one modification resolves 3-4 test failures by targeting root causes) from sequential trial-and-error (ratio 1.07: one modification per failure). Controlled experiments on mock agents and synthetic datasets validate a two-stage mechanism with large effect sizes: agents cluster errors by shared characteristics at $2\times$ random rate (clustering coefficient 1.909 vs 0.950, p=0.001), then prioritize high-impact fixes yielding 2.$15\times$ higher proportion of multi-test modifications (42.5\% vs 19.8\%). However, pattern-based transfer to held-out tests failed (slope ratio 0.82 < threshold 1.5, p=0.504, 0\% pattern usage), revealing strategic debugging operates as a feedback loop effective within revealed test execution, not a learning system enabling autonomous prediction without test execution. The framework establishes feasibility of process-based debugging metrics on controlled mock agents; ecological validity with production systems remains unverified.
